%%==========================================================
%%  RIESGOS ÉTICOS
%%  Se incluye como sección final del Capítulo 1.
%%==========================================================

\section{Riesgos éticos}
\label{sec:riesgos-eticos}

Un sistema que estandariza y automatiza el inicio de aplicaciones concentra
en un único punto decisiones técnicas que después se replican en varios
proyectos. En un contexto financiero esa concentración importa más, por el
manejo de información sensible y de lógica de negocio regulada. El proyecto
identificó cuatro riesgos éticos al plantearse, cada uno con la mitigación a
la que se comprometía, y lo que corresponde decir aquí es cuál de las cuatro
se cumplió. El veredicto final de cada uno, con la evidencia de la
evaluación, lo da el Capítulo~\ref{cap:evaluacion}.

\paragraph*{Exposición de información sensible en plantillas y módulos}

El compromiso era separar de forma estricta la estructura base de los valores
de configuración sensibles, promoviendo variables de entorno y mecanismos de
gestión de secretos desde el diseño. \textbf{Se cumplió, y quedó verificado
por ejecución.} Es lo que la especificación de requerimientos exige como
atributo de configurabilidad y lo que el artefacto hace, porque ninguna
dirección ni credencial vive en el código del producto, toda
configuración entra por variable de entorno validada contra un esquema, el
circuito de verificación barre el código del producto derivado buscando
valores fijos y direcciones, y el registro de eventos del producto redacta
por nombre de clave todo secreto o identificador de persona. La verificación
automática es más de lo que el compromiso pedía.

\paragraph*{Propagación de lógica financiera mal interpretada}

El compromiso era someter a revisión técnica especializada los módulos con
lógica financiera antes de incorporarlos al catálogo. \textbf{El riesgo no se
materializó}, porque ninguno de los dos módulos construidos lleva lógica
financiera, ya que uno guarda y lee registros y el otro exporta señales de
observabilidad. Las dos capacidades son técnicas y transversales por decisión
del marco, que reparte lo técnico al módulo y el dominio al producto. Lo que
queda escrito es la condición de reactivación, porque el día que el catálogo
reciba un módulo que interprete un estándar financiero, y el de mensajería
excluido en la sección~\ref{sec:fundamentos-adoptados} es el caso obvio, la
revisión especializada es exigible antes de su entrada.

\paragraph*{Exclusión de equipos con menor capacidad técnica}

El compromiso tenía dos partes y \textbf{se cumplió a medias}. La primera era
elaborar documentación accesible orientada a distintos niveles de
experiencia, y \textbf{no se construyó}. La propuesta existió, y se descartó
con un argumento de oportunidad, que esa documentación crece con los módulos
incorporados y hacerla antes documenta un producto que todavía va a cambiar
de forma. El argumento es razonable como decisión de alcance de una fase y no
es una respuesta al riesgo ético, porque nadie advirtió que lo fuera, ya que
el rechazo comparó la propuesta contra el calendario y no contra el
compromiso. Se declara así, y la documentación por niveles queda como
recomendación de adopción.

La segunda parte sí está. El producto derivado nace con documentación que
explica qué es, qué hace cada operación, cuál lleva efecto de negocio y cuál
es un andamio; el recorrido de creación se lee en lenguaje de negocio con el
nombre del estándar al lado; la herramienta declara su alcance antes de
escribir y pide confirmación; y los mensajes de error dicen qué encontró, qué
esperaba y qué resuelve la situación. Todo eso baja la barrera de adopción,
que es lo que el riesgo persigue. El acompañamiento con procesos de formación
y transferencia de conocimiento es organizacional y queda fuera de lo que un
artefacto puede resolver.

\paragraph*{Concentración de poder técnico en quienes controlan el sistema}

El compromiso era definir un modelo de gobernanza distribuido con procesos
participativos, conocimiento documentado de forma exhaustiva y una evolución
que responda a necesidades colectivas. \textbf{El mecanismo existe y el
ejercicio no}, y la razón es una condición del entorno declarada desde la
segunda fase.

El mecanismo está completo, porque el catálogo declara un responsable por
activo, la aprobación del gobierno de la línea es obligatoria sobre la ruta
del manifiesto, y el procedimiento de incorporación tiene sus pasos y sus
condiciones de validez verificables por tipo de activo. La documentación
exhaustiva es la parte más fuerte del compromiso, con casi doscientas
decisiones registradas con su contexto, su alternativa y su traza, más de
treinta documentos de expediente y un marco con su propio mecanismo de
evolución, que es lo que permite que otra persona retome el sistema. Y la
evolución respondió a necesidades que no son del autor, porque las doce
necesidades y las tres brechas salen de las entrevistas con las personas del
contexto, y varias rondas de mejora las abrieron una revisión externa y
reportes de uso.

Lo que no hay es más de un participante. El proyecto se construyó con una
sola persona en los tres roles, de modo que proponente y revisor coinciden, y
una gobernanza distribuida instanciada con un participante es un mecanismo
probado y un ejercicio ausente. Trasladarlo a operación consiste en cambiar
quién figura como responsable de cada activo, lo que el modelo de gobierno ya
contempla, y queda como recomendación de adopción.
